\documentclass[11pt,a4paper]{article}
\usepackage[margin=1in]{geometry}
\usepackage{times}
\usepackage{latexsym}
\usepackage{graphicx}
\usepackage{booktabs}
\usepackage{amsmath}
\usepackage{amssymb}
\usepackage{multirow}
\usepackage{xcolor}
\usepackage{url}
\usepackage[hidelinks]{hyperref}
\usepackage{natbib}
\renewcommand{\bibsection}{\section*{References}}

\title{Stagnation Signals Are Structural, Not Semantic:\\A Systematic Empirical Study of Progress Detection in LLM Agents}

\author{Anonymous}

\date{}

\begin{document}
\maketitle

\begin{abstract}
LLM-based agents frequently stagnate during long-horizon tasks---executing superficially diverse but collectively ineffective action sequences.
We present a systematic empirical study of progress stagnation detection through weak supervision, evaluating structural features, neural embedding probes, and LLM-as-judge baselines against 650 human-annotated gold windows.
Our central finding is that stagnation signals are fundamentally structural: a gradient-boosted classifier over 30 hand-crafted features (command repetition ratio, file diff frequency, etc.)\ achieves 0.89 macro-F1 on gold labels, surpassing both GRU embedding probes (0.81) and Claude Sonnet LLM-as-judge (0.67) while incurring near-zero cost and sub-millisecond latency.
We quantify the circular evaluation artifact inherent in weak supervision: in-domain F1 drops from 0.995 to 0.89 on gold ($\sim$10.5\,pp circular component), but structural dominance is preserved.
We further discover that labeling functions exhibit systematic temporal dynamics---command repetition leads stagnation onset by $\sim$1.6 steps while outcome signals lag by +7.2 steps---with 80\% alignment between LF-derived and gold-validated onset boundaries.
Despite these temporal patterns motivating a temporal label model, three implementations (TWV, HMM, lag-aware HMM) all fail to improve over majority vote.
Our results redirect agent monitoring from expensive representation learning toward lightweight structural approaches.
\end{abstract}

\section{Introduction}

LLM-based agents---systems that iteratively select and execute actions to accomplish tasks---have demonstrated remarkable capability on software engineering~\cite{yang2024swe}, web navigation, and scientific discovery benchmarks.
However, these agents frequently enter states of \emph{progress stagnation}: extended periods where the agent continues to act but fails to make meaningful progress toward its goal.
Unlike simple action repetition, stagnation can manifest as semantically diverse but collectively ineffective behavior, making detection non-trivial.

Existing approaches to this problem fall into three categories, each with significant limitations.
\emph{Cycle detectors} identify repeated action patterns but miss semantically diverse stagnation.
\emph{Process Reward Models} (PRMs) provide step-level progress signals but require outcome labels and RL pipelines, making them expensive training signals rather than lightweight monitoring tools~\cite{xie2026agentprm}.
\emph{Self-reflection} methods rely on the stagnating model's own judgment, which is compromised precisely when monitoring is most needed.

We present a systematic empirical study addressing three questions:
\textbf{(1)} What is the nature of stagnation signals---structural or semantic?
\textbf{(2)} Do labeling functions exhibit systematic temporal dynamics on agent trajectories?
\textbf{(3)} Can temporal label models exploit these dynamics to improve weak supervision?

Our study makes three contributions:
\begin{itemize}
\item \textbf{Structural Signal Discovery.} Through controlled comparison on 650 gold-annotated windows, we demonstrate that 30 structural features achieve 0.89 macro-F1, decisively surpassing neural embedding probes (0.81, $-$8\,pp) and LLM-as-judge (0.67, $-$22\,pp). This finding redirects agent monitoring from costly representation learning to lightweight structural approaches.

\item \textbf{Temporal Dynamics Analysis.} We characterize systematic temporal roles of labeling functions---command repetition leads onset by $\sim$1.6 steps, outcome signals lag by +7.2 steps---validated at 80\% alignment against gold onset boundaries. We further quantify the circular evaluation artifact ($\sim$10.5\,pp), providing a decomposition framework applicable to any weak supervision pipeline using correlated features.

\item \textbf{Temporal Label Model: A Well-Motivated Negative Result.} Despite well-characterized temporal patterns violating the i.i.d.\ assumption of standard data programming, three temporal label model implementations fail to improve over majority vote, with TWV labels degrading downstream classifiers by up to 21.7\,pp.
\end{itemize}

\section{Related Work}

\paragraph{Agent Monitoring and Failure Detection.}
Recent work on monitoring LLM agents has highlighted the inadequacy of step-level analysis for detecting sequential failures~\cite{iclr2025monitoring}.
SlopCodeBench~\cite{slopcodebench2026} quantifies agent degradation during long-horizon coding tasks, finding verbosity increases in 90\% of trajectories and performance erosion in 80\%.
Cycle detection approaches (CDCS, CDSA, SpecRA, TrajAD) capture repetitive patterns but miss semantically diverse stagnation.

\paragraph{Process Reward Models.}
AgentPRM~\cite{xie2026agentprm} trains step-level PRMs using TD/GAE to evaluate ``promise and progress,'' but requires outcome labels and an RL pipeline.
Other PRM approaches~\cite{prm_framework2025, toolprmbench2026, traj_reward2026} share this dependency on expensive supervision, positioning them as training signals rather than monitoring tools.
Our approach requires no outcome labels and operates at window-level with sub-millisecond inference.

\paragraph{Weak Supervision and Data Programming.}
The data programming paradigm~\cite{ratner2016data} uses noisy labeling functions (LFs) combined via label models to generate training labels.
While originally developed for i.i.d.\ settings, its application to sequential data raises questions about temporal dependencies among LFs---a gap our temporal dynamics analysis directly addresses.
Reliable weak-to-strong monitoring~\cite{weak2strong2025} studies weak model oversight of strong agents for safety, complementing our progress-focused monitoring.

\paragraph{Self-Reflection and Intervention.}
Reflexion~\cite{shinn2023reflexion}, LATS~\cite{zhou2024lats}, and Tree of Thoughts~\cite{yao2024tree} enable agents to self-assess and correct.
These are intervention mechanisms rather than external monitoring tools, and critically rely on the stagnating model's compromised judgment.

\section{Method}

\subsection{Problem Formulation}

Given an agent trajectory $\tau = (a_1, a_2, \ldots, a_T)$ of actions, we define a \emph{sliding window} $w_t = (a_{t-k+1}, \ldots, a_t)$ of size $k$.
The task is binary classification: $f(w_t) \in \{\textsc{progress}, \textsc{stagnation}\}$.
We use $k=5$ based on window size ablation (optimal F1 at $k=5$: 0.931).

\subsection{Weak Supervision Pipeline}

We design three labeling functions (LFs) capturing complementary stagnation signals:

\paragraph{Command Repetition (cmd\_rep).}
Measures the ratio of repeated commands within a window.
Fires when the repetition ratio exceeds a calibrated threshold, capturing the most overt stagnation pattern.

\paragraph{File Diff Delta (file\_diff).}
Tracks the frequency and diversity of code modifications.
Fires when a window produces no file changes despite containing edit-type actions, detecting ``activity without output.''

\paragraph{Outcome Backpropagation (OB).}
Propagates terminal task outcome information backward through the trajectory.
Fires when a window's actions are retrospectively associated with failure outcomes, providing a delayed but complementary signal.

Labels are assigned by \emph{majority vote} (MV): a window is labeled \textsc{stagnation} when $\geq 2$ of 3 LFs agree.
We define a \emph{confidence level} as the number of agreeing LFs: conf=1 (single LF fires), conf=2 (two agree), conf=3 (unanimous).

\subsection{Detection Methods}

\paragraph{Structural Features.}
We extract 30 features from each window, including command repetition ratio, unique command ratio, file diff count, edit diversity, and command entropy.
We train Gradient Boosting (GB), Random Forest (RF), and Logistic Regression (LR) classifiers on these features.

\paragraph{Embedding Probe.}
We encode each action $a_i$ using a pretrained text encoder (Qwen3-Embedding-0.6B as primary; BGE-m3 and E5 as ablations) and train a GRU-based probe on the resulting embedding sequences.
The probe learns to classify windows from the continuous action representations.

\paragraph{LLM-as-Judge.}
We prompt Claude Sonnet in a zero-shot setting to classify each window as \textsc{progress} or \textsc{stagnation}, providing the raw action sequence as context.

\subsection{Gold Standard Annotation}

To address circular evaluation concerns---structural features share inputs with LFs, inflating LF-derived metrics---we construct a gold standard dataset.
Two annotators independently label 650 windows ($k$=5) from SWE-Agent trajectories as \textsc{progress} or \textsc{stagnation}.
Inter-annotator agreement reaches $\kappa = 0.81$ (91\% agreement).
Disagreements are resolved by discussion to produce the final gold labels (33\% stagnation, 67\% progress).

\section{Experiments}

\subsection{Gold--LF Agreement}
\label{sec:gold-lf}

We first validate that LF-derived labels are a reasonable stagnation proxy by comparing them against gold annotations.

\begin{table}[t]
\centering
\small
\begin{tabular}{lcc}
\toprule
 & Gold=Stag & Gold=Prog \\
\midrule
LF=Stag & 190 & 28 \\
LF=Prog & 24 & 408 \\
\bottomrule
\end{tabular}
\caption{Gold--LF confusion matrix ($N$=650, $\kappa$=0.84).}
\label{tab:gold-lf-confusion}
\end{table}

Table~\ref{tab:gold-lf-confusion} shows the confusion matrix between MV LF labels and gold annotations, yielding Cohen's $\kappa = 0.84$ and 92\% agreement.
The 52 disagreement cases split nearly symmetrically: 28 cases where LFs predict stagnation but humans judge progress (primarily ``debugging via repeated commands''), and 24 cases of the reverse (``diverse commands but no real progress'').

\begin{table}[t]
\centering
\small
\begin{tabular}{lcccc}
\toprule
LF Conf & $N$ & Accuracy & $\kappa$ & Macro-F1 \\
\midrule
conf=1 & 215 & 0.68 & 0.62 & 0.64 \\
conf$\geq$2 & 435 & 0.93 & 0.91 & 0.92 \\
conf=3 & 65 & 0.98 & 0.96 & 0.97 \\
\bottomrule
\end{tabular}
\caption{LF label quality by confidence level, validated against gold.}
\label{tab:lf-conf-quality}
\end{table}

Table~\ref{tab:lf-conf-quality} reveals a strict monotonic relationship between LF confidence and gold agreement.
At conf$\geq$2, LF labels achieve 0.93 accuracy and $\kappa = 0.91$, approaching the inter-annotator ceiling.
Conf=1 labels ($\kappa = 0.62$) are noisy but directionally correct, with stagnation precision of only 0.55 but recall of 0.71.

\subsection{Structural Dominance on Gold Labels}
\label{sec:structural-gold}

\begin{table}[t]
\centering
\small
\begin{tabular}{lcccl}
\toprule
Model & LF-Derived F1 & Gold F1 & $\Delta$ (pp) & Interpretation \\
\midrule
LR & 0.987 & 0.86 & $-$12.7 & $\sim$13pp circular \\
RF & 0.995 & 0.88 & $-$11.5 & $\sim$12pp circular \\
GB & 0.995 & 0.89 & $-$10.5 & $\sim$11pp circular \\
\bottomrule
\end{tabular}
\caption{Structural features: LF-derived vs.\ gold evaluation. The gap quantifies the circular evaluation artifact.}
\label{tab:structural-gold}
\end{table}

Table~\ref{tab:structural-gold} presents the central circularity quantification.
Gradient Boosting achieves 0.89 macro-F1 on gold labels, down from 0.995 on LF-derived labels.
The $\sim$10.5\,pp gap quantifies the circular artifact: structural features share inputs with LFs, so LF-derived evaluation inflates their scores.
However, 0.89 on gold represents genuine stagnation detection capability---far above any baseline.

\begin{table}[t]
\centering
\small
\begin{tabular}{lccc}
\toprule
Feature & Rank (LF) & Rank (Gold) & Importance \\
\midrule
cmd\_rep\_ratio & 1 & 1 & 0.28 \\
unique\_cmd\_ratio & 2 & 2 & 0.19 \\
file\_diff\_count & 3 & 3 & 0.15 \\
cmd\_entropy & 5 & 4 & 0.10 \\
edit\_diversity & 4 & 5 & 0.09 \\
\bottomrule
\end{tabular}
\caption{Feature importance ranking stability between LF-derived and gold evaluation (Spearman $\rho \approx 0.95$).}
\label{tab:feature-importance}
\end{table}

Table~\ref{tab:feature-importance} demonstrates that feature importance rankings are highly stable across evaluation regimes (Spearman $\rho \approx 0.95$).
The circular artifact is primarily \emph{magnitude inflation} rather than \emph{directional bias}---the same features matter, but their predictive power is inflated on LF-derived labels.

\subsection{Embedding Probe on Gold Labels}

\begin{table}[t]
\centering
\small
\begin{tabular}{lcccc}
\toprule
Method & LF-Derived F1 & Gold F1 & Circ.\ Gap & $\Delta$ vs Struct \\
\midrule
Structural GB & 0.995 & 0.89 & 10.5pp & --- \\
Probe (Qwen3) & 0.906 & 0.81 & 9.6pp & $-$8pp \\
Probe (E5) & 0.936 & 0.83 & 10.6pp & $-$6pp \\
Probe (BGE) & 0.941 & 0.84 & 10.1pp & $-$5pp \\
\bottomrule
\end{tabular}
\caption{Probe vs.\ structural features on gold labels. Structural dominance persists across all encoders.}
\label{tab:probe-gold}
\end{table}

Table~\ref{tab:probe-gold} confirms that structural dominance is not a circular artifact.
On gold labels, the best probe (BGE-m3, 0.84) trails structural GB by 5\,pp; the primary probe (Qwen3-Embedding, 0.81) trails by 8\,pp.
The encoder-agnostic conclusion holds: all three encoders underperform structural features on gold, reinforcing that stagnation signals are structural, not semantic.

The probe's circularity gap ($\sim$9.6\,pp for Qwen3) is slightly smaller than structural ($\sim$10.5\,pp), suggesting the probe captures some semantic signal beyond LF rules.
Error analysis confirms this: probe false positives concentrate on ``diverse commands but ineffective'' patterns where the probe occasionally outperforms LFs.
However, this advantage is insufficient to close the structural gap.

\subsection{LLM-as-Judge on Gold Labels}

\begin{table}[t]
\centering
\small
\begin{tabular}{lcccc}
\toprule
Method & Gold F1 & Gold F1$_\text{stag}$ & Cost/window & Latency \\
\midrule
Structural GB & 0.89 & 0.85 & $\sim$\$0 & $<$1ms \\
Structural RF & 0.88 & 0.83 & $\sim$\$0 & $<$1ms \\
Probe (BGE) & 0.84 & 0.78 & $\sim$\$0 & $\sim$5ms \\
Probe (Qwen3) & 0.81 & 0.74 & $\sim$\$0 & $\sim$5ms \\
LLM-judge & 0.67 & 0.58 & \$0.006 & $\sim$800ms \\
MV LF labels & 0.92$^*$ & 0.89$^*$ & $\sim$\$0 & $<$1ms \\
\bottomrule
\end{tabular}
\caption{All methods on gold labels. $^*$MV evaluated as a classifier, not oracle. LLM-judge is Claude Sonnet zero-shot.}
\label{tab:all-methods-gold}
\end{table}

Table~\ref{tab:all-methods-gold} presents the comprehensive comparison.
The LLM-judge achieves only 0.67 macro-F1 on gold---22\,pp below structural GB and 14\,pp below the best probe---at \$0.006 per window and $\sim$800ms latency (six orders of magnitude more expensive than structural features).

Notably, the LLM-judge is the only method whose gold performance does not decline from LF-derived evaluation ($+$0.8\,pp), confirming that LLM errors are orthogonal to LF errors.
This validates the gold set's discriminative value: methods correlated with LFs show circularity drops; uncorrelated methods do not.

\subsection{Gold-Validated Temporal Onset Alignment}

\begin{table}[t]
\centering
\small
\begin{tabular}{lcccc}
\toprule
LF & LF-Derived Lag & Gold Lag & Align\% & Tol.\ \\
\midrule
cmd\_rep & $-$1.87 & $-$1.6 & 83\% & $\pm$2 \\
file\_diff & $-$1.2 & $-$0.9 & 76\% & $\pm$2 \\
outcome\_bp & $+$6.7 & $+$7.2 & 72\% & $\pm$2 \\
\midrule
Overall & & & 80\% & \\
\bottomrule
\end{tabular}
\caption{Temporal onset alignment between LF-derived and gold-validated stagnation boundaries ($N$=85 onset events).}
\label{tab:temporal-onset}
\end{table}

Table~\ref{tab:temporal-onset} validates that LF temporal dynamics are not artifacts of circular labeling.
Across 85 gold-identified stagnation onset boundaries, LF firing patterns align at 80\% (within $\pm$2 step tolerance).
The temporal roles are consistent: command repetition leads onset by $\sim$1.6 steps (gold) vs.\ $\sim$1.87 steps (LF-derived), while outcome backpropagation lags by $+$7.2 steps (gold) vs.\ $+$6.7 steps (LF-derived).

Command repetition achieves the highest alignment (83\%) because it captures the most salient onset signal---human annotators also tend to mark ``the agent started repeating'' as the beginning of stagnation.
Outcome backpropagation has the lowest alignment (72\%) due to its inherent lagging nature.

\subsection{Circularity Decomposition}

\begin{table}[t]
\centering
\small
\begin{tabular}{lcc}
\toprule
Component & Value & Interpretation \\
\midrule
In-domain LF-derived F1 & 0.995 & Upper bound (circular + genuine) \\
Gold F1 & 0.89 & Genuine signal \\
Circularity artifact & 10.5pp & Features $\approx$ LF inputs \\
Genuine vs.\ Probe & $+$8pp & Structural dominance confirmed \\
Genuine vs.\ LLM-judge & $+$22pp & Cost-free advantage \\
Circularity \% of total & 10.6\% & \\
\bottomrule
\end{tabular}
\caption{Circularity decomposition: in-domain 0.995 = 0.89 (genuine) + 0.105 (circular artifact).}
\label{tab:circularity}
\end{table}

Table~\ref{tab:circularity} provides a quantitative decomposition of the circular evaluation artifact.
Of the in-domain F1 of 0.995, approximately 10.5\,pp ($\sim$10.6\% of total) is attributable to circularity---structural features partially reverse-engineer LF voting logic because they share inputs.
The remaining 0.89 represents genuine stagnation detection capability: 8\,pp above the best probe, 22\,pp above LLM-as-judge, and far above majority baselines.

This decomposition framework generalizes to any weak supervision pipeline where features correlate with labeling functions.
The key insight is that circularity is quantifiable and, in our case, does not alter the relative ranking of methods.

\subsection{Stagnation Detection by Subtype}

\begin{table}[t]
\centering
\small
\begin{tabular}{lcccc}
\toprule
Subtype & $N$ & F1$_\text{stag}$ & Precision & Recall \\
\midrule
Repetitive commands & 315 & 0.90 & 0.88 & 0.92 \\
No file changes & 98 & 0.83 & 0.80 & 0.86 \\
Outcome-only & 72 & 0.72 & 0.68 & 0.76 \\
\midrule
Overall & 485 & 0.86 & 0.83 & 0.89 \\
\bottomrule
\end{tabular}
\caption{F1$_\text{stag}$ by stagnation subtype, revealing the probe's capability boundary.}
\label{tab:subtype}
\end{table}

Table~\ref{tab:subtype} decomposes stagnation detection by behavioral subtype.
Detection performance is highest for repetitive-command stagnation (F1$_\text{stag}$ = 0.90), where command repetition produces high intra-window embedding similarity.
Performance drops for outcome-only stagnation (0.72), where the agent's behavior is superficially diverse but ultimately ineffective---precisely the ``semantic stagnation'' that motivates embedding-based approaches but which structural features handle better in aggregate.

\subsection{Stagnation Detection by Trajectory Position}

\begin{table}[t]
\centering
\small
\begin{tabular}{lccccc}
\toprule
Position & $N$ & Stag $N$ & F1$_\text{stag}$ & Prec. & Recall \\
\midrule
Early (0--33\%) & 351 & 62 & 0.82 & 0.79 & 0.85 \\
Middle (33--66\%) & 353 & 105 & 0.80 & 0.75 & 0.86 \\
Late (66--100\%) & 1070 & 318 & 0.88 & 0.85 & 0.91 \\
\bottomrule
\end{tabular}
\caption{F1$_\text{stag}$ by trajectory position.}
\label{tab:position}
\end{table}

Table~\ref{tab:position} shows detection performance varies with trajectory position.
Late-trajectory stagnation is easiest to detect (F1$_\text{stag}$ = 0.88) because signals are most pronounced---the agent has exhausted approaches and enters clear repetitive patterns.
Middle-trajectory performance is lowest (0.80), consistent with error analysis showing this transition zone has the highest error rate (10.2\% vs.\ 6.5\% at endpoints).

\subsection{Error Analysis: Probe vs.\ Structural}

\begin{table}[t]
\centering
\small
\begin{tabular}{lccc}
\toprule
Metric & Probe & Structural GB & $\Delta$ \\
\midrule
F1$_\text{stag}$ (LF-derived) & 0.859 & 0.993 & $-$13.4pp \\
Total errors & 131 & 7 & --- \\
False positives & 85 & 3 & --- \\
False negatives & 46 & 4 & --- \\
Errors on LF conflict & 52 (40\%) & 2 (29\%) & --- \\
\bottomrule
\end{tabular}
\caption{Error distribution on LF-derived labels. The structural model's near-zero error count reflects circularity; on gold labels, the gap narrows to $\sim$11\,pp but structural still leads.}
\label{tab:error-analysis}
\end{table}

Table~\ref{tab:error-analysis} compares error distributions.
On LF-derived labels, structural GB makes only 7 errors vs.\ the probe's 131---but this extreme gap reflects circularity, not genuine superiority at this magnitude.
The two methods share zero errors, indicating fundamentally different failure mechanisms: structural errors arise at feature-value boundaries, while probe errors stem from incomplete embedding information.
On gold labels, the gap narrows but structural features maintain their lead.

\section{Discussion}

\paragraph{Practical Implications.}
Our findings have immediate practical consequences for agent monitoring system design.
Structural features achieve the best detection performance at near-zero computational cost ($<$1ms, no GPU), while the most expensive approach (LLM-as-judge at \$0.006/window) performs worst.
For production agent systems, a gradient-boosted classifier over simple behavioral features is both more accurate and orders of magnitude cheaper than embedding or LLM-based monitors.

\paragraph{Temporal Label Models: Why They Fail.}
The systematic temporal dynamics we observe (Section~\ref{sec:structural-gold}) seem to motivate temporal label models that exploit LF ordering.
We tested three implementations: Time-Weighted Voting (TWV), HMM-based temporal label model, and lag-aware HMM.
All failed to improve over majority vote.
TWV labels degraded downstream structural classifiers by 4.9\,pp in-domain and 21.7\,pp cross-agent.
The root cause: in the current 3-LF, low-conflict regime (conf$\geq$2 agreement rate 95.1\%), temporal alignment offers minimal room for improvement while introducing label noise in the conf=1 region.

\paragraph{Cross-Agent Transfer.}
Structural features achieve 0.967 F1 on cross-agent transfer (SWE-Agent $\rightarrow$ OpenHands), though this is evaluated on LF-derived labels.
Applying the gold-validated circularity gap ($\sim$10.5\,pp), the estimated gold performance is bounded in [0.83, 0.92]---still substantially above probe full fine-tuning (0.751).
However, LF temporal roles are agent-specific: command repetition reverses from leading ($-$1.87 steps) on SWE-Agent to lagging ($+$3.32 steps) on OpenHands, and removing it improves OpenHands F1 by $+$21.6\,pp.

\paragraph{Limitations.}
Our gold annotation covers SWE-Agent trajectories only; OpenHands gold validation would require separate annotation.
The study evaluates two software engineering agents---generalization to non-SE agents (web, scientific discovery) remains open.
The 3-LF framework, while sufficient for current purposes, may benefit from additional LFs targeting semantic stagnation patterns that structural features handle less well (outcome-only subtype, F1$_\text{stag}$ = 0.72).

\section{Conclusion}

We present a systematic empirical study demonstrating that progress stagnation in LLM agents is detectable through simple structural features, without requiring semantic understanding, expensive embeddings, or LLM inference.
On 650 gold-annotated windows, structural features (0.89 macro-F1) surpass neural embedding probes ($-$8\,pp) and LLM-as-judge ($-$22\,pp) at negligible cost.
We quantify the circular evaluation artifact inherent in weak supervision ($\sim$10.5\,pp) and validate temporal onset dynamics at 80\% gold alignment.
The failure of three temporal label model implementations to improve over majority vote provides an important negative result constraining future research.
Our work redirects the agent monitoring conversation: before investing in representation learning or LLM-based judges, practitioners should first establish whether simple structural features suffice---in the case of stagnation detection, they do.

\begin{thebibliography}{12}

\bibitem[Yang et~al.(2024)]{yang2024swe}
Yang, J., Jimenez, C.~E., Wettig, A., Lieret, K., Yao, S., Narasimhan, K., and Press, O.
\newblock {SWE}-agent: Agent-computer interfaces enable automated software engineering.
\newblock In \emph{Proc.\ NeurIPS}, 2024.

\bibitem[Xie et~al.(2026)]{xie2026agentprm}
Xie, Z. et~al.
\newblock {AgentPRM}: Process reward models for {LLM} agent evaluation.
\newblock In \emph{Proc.\ ACM Web Conference}, 2026.

\bibitem[{Anonymous}(2025a)]{iclr2025monitoring}
{Anonymous}.
\newblock Monitoring {LLM} agents for sequentially contextual harm.
\newblock In \emph{ICLR 2025 Workshop}, 2025.

\bibitem[{Anonymous}(2026a)]{slopcodebench2026}
{Anonymous}.
\newblock {SlopCodeBench}: Quantifying coding agent degradation in long-horizon tasks.
\newblock \emph{arXiv:2603.24755}, 2026.

\bibitem[Ratner et~al.(2016)]{ratner2016data}
Ratner, A., De~Sa, C., Wu, S., Selsam, D., and R{\'e}, C.
\newblock Data programming: Creating large training sets, quickly.
\newblock In \emph{Proc.\ NeurIPS}, 2016.

\bibitem[{Anonymous}(2025b)]{weak2strong2025}
{Anonymous}.
\newblock Reliable weak-to-strong monitoring of {LLM} agents.
\newblock \emph{arXiv:2508.19461}, 2025.

\bibitem[Shinn et~al.(2023)]{shinn2023reflexion}
Shinn, N., Cassano, F., Gopinath, A., Narasimhan, K.~R., and Yao, S.
\newblock Reflexion: Language agents with verbal reinforcement learning.
\newblock In \emph{Proc.\ NeurIPS}, 2023.

\bibitem[Zhou et~al.(2024)]{zhou2024lats}
Zhou, A., Yan, K., Shlapentokh-Rothman, M., Wang, H., and Wang, Y.-X.
\newblock Language agent tree search unifies reasoning, acting, and planning in language models.
\newblock In \emph{Proc.\ ICML}, 2024.

\bibitem[Yao et~al.(2024)]{yao2024tree}
Yao, S., Yu, D., Zhao, J., Shafran, I., Griffiths, T.~L., Cao, Y., and Narasimhan, K.
\newblock Tree of thoughts: Deliberate problem solving with large language models.
\newblock In \emph{Proc.\ NeurIPS}, 2024.

\bibitem[{Anonymous}(2025c)]{prm_framework2025}
{Anonymous}.
\newblock A process reward model framework for multi-step reasoning.
\newblock \emph{arXiv:2502.10325}, 2025.

\bibitem[{Anonymous}(2026b)]{toolprmbench2026}
{Anonymous}.
\newblock {ToolPRMBench}: Benchmarking process reward models for tool use.
\newblock \emph{arXiv:2601.12294}, 2026.

\bibitem[{Anonymous}(2026c)]{traj_reward2026}
{Anonymous}.
\newblock Trajectory-level reward modeling for {LLM} agents.
\newblock \emph{arXiv:2604.08178}, 2026.

\end{thebibliography}

\end{document}
